\documentclass[11pt, a4paper]{article}
\usepackage{amsmath, amssymb, booktabs, geometry, hyperref, xcolor, listings}
\geometry{margin=2.5cm}
\hypersetup{colorlinks=true, linkcolor=blue}
\lstset{basicstyle=\ttfamily\small, backgroundcolor=\color{gray!10}, frame=single}

\title{\textbf{GJR-GARCH(1,1,1) Model}\\[4pt]
\large Portfolio VaR Engine --- Model Documentation}
\author{Risk Modelling --- Trade Republic}
\date{April 2026}

\begin{document}
\maketitle

%-----------------------------------------------------------------------
\section{Overview}

The GJR-GARCH model (Glosten, Jagannathan, and Runkle, 1993) extends the
standard GARCH(1,1) specification with an asymmetric term that allows
negative return shocks to have a larger impact on future volatility than
positive shocks of the same magnitude --- the so-called \textit{leverage
effect}. This is especially relevant for equity portfolios.

\textbf{Engine key:} \texttt{'gjr'} \\
\textbf{Sources:}
\begin{itemize}
  \item \texttt{risk\_models/garch\_models.py}, lines 38--54 (\texttt{garch\_gjr})
  \item \texttt{risk\_models/garch\_models.py}, lines 58--131 (\texttt{calculate\_vaR\_garch})
  \item \texttt{risk\_analytics/risk\_engines.py}, lines 135--152 (engine dispatch)
\end{itemize}
\textbf{Library:} \texttt{arch} (Kevin Sheppard)

%-----------------------------------------------------------------------
\section{Volatility Model}

\subsection{GJR-GARCH(1,1,1) Specification}

Let $r_t$ denote the portfolio return at time $t$. The model assumes:
\begin{equation}
r_t = \mu + \varepsilon_t, \qquad \varepsilon_t = \sigma_t z_t,
\qquad z_t \overset{i.i.d.}{\sim} D(0,1)
\end{equation}
where $D(0,1)$ is a standardised distribution (Normal or Student-$t$),
and $\sigma_t^2$ follows:
\begin{equation}
\sigma_t^2 = \omega
  + \alpha \, \varepsilon_{t-1}^2
  + \gamma \, \varepsilon_{t-1}^2 \, \mathbb{1}_{\{\varepsilon_{t-1} < 0\}}
  + \beta \, \sigma_{t-1}^2
\label{eq:gjr}
\end{equation}

\begin{itemize}
  \item $\omega > 0$: long-run variance intercept
  \item $\alpha \geq 0$: ARCH effect (symmetric impact of shocks)
  \item $\gamma \geq 0$: asymmetry coefficient (leverage effect;
        positive shocks have impact $\alpha$, negative shocks have
        impact $\alpha + \gamma$)
  \item $\beta \geq 0$: GARCH persistence
  \item Stationarity requires: $\alpha + \tfrac{1}{2}\gamma + \beta < 1$
\end{itemize}

\subsection{ARCH library call}

The model is fitted via:
\begin{lstlisting}[language=Python]
arch_model(rets_rescaled, p=1, q=1, o=1, vol='GARCH', dist=dist)
\end{lstlisting}
where \texttt{o=1} adds the single asymmetric (GJR) term,
\texttt{p=1} is the GARCH lag order, and \texttt{q=1} is the ARCH lag order.
Returns are rescaled by a factor (default 1000) before fitting for
numerical stability.

\subsection{EWMA Pre-weighting (optional)}

Before fitting, an optional EWMA decay can be applied to down-weight
older observations:
\begin{equation}
\tilde{r}_t = \text{EWMA}_\lambda(r_t), \qquad \lambda \in (0,1]
\end{equation}
With $\lambda = 1$ (default for GARCH engines) no decay is applied.

%-----------------------------------------------------------------------
\section{VaR and CVaR Calculation}

\subsection{One-Step Forecast}

After fitting, a one-step-ahead forecast of conditional mean and variance
is generated using \texttt{ARCHModelResult.forecast()}:
\begin{equation}
\hat{\mu}_{t+1} = \hat{\mu}, \qquad \hat{\sigma}_{t+1}^2 = \text{forecast variance}
\end{equation}
The forecast uses 1000 simulations internally for the variance path
(\texttt{simulations=1000}).

\subsection{Parametric VaR (without FHS)}

\textbf{Value-at-Risk:}
\begin{equation}
\widehat{\text{VaR}}_\alpha
= -\hat{\mu}_{t+1} + \hat{\sigma}_{t+1} \cdot q_\alpha^D
\label{eq:gjr_var}
\end{equation}
where $q_\alpha^D$ is the $\alpha$-quantile of the fitted distribution $D$.

For Normal innovations: $q_\alpha^N = \Phi^{-1}(\alpha) < 0$. \\
For Student-$t$ innovations: $q_\alpha^t = t^{-1}(\alpha; \hat{\nu}) < 0$.

\textbf{CVaR --- Normal distribution:}
\begin{equation}
\widehat{\text{CVaR}}_\alpha^{\text{N}}
= -\hat{\mu}_{t+1} + \hat{\sigma}_{t+1} \cdot
  \frac{\phi\!\left(\Phi^{-1}(\alpha)\right)}{\alpha}
\label{eq:gjr_cvar_norm}
\end{equation}

\textbf{CVaR --- Student-$t$ distribution:}
\begin{equation}
\widehat{\text{CVaR}}_\alpha^{t}
= -\hat{\mu}_{t+1}
  - \frac{\hat{\sigma}_{t+1}}{\alpha(1 - \hat{\nu})}
    \Bigl(\hat{\nu} - 2 + \bigl[t^{-1}(\alpha;\hat{\nu})\bigr]^2\Bigr)
    \cdot f_t\!\bigl(t^{-1}(\alpha;\hat{\nu});\hat{\nu}\bigr)
\label{eq:gjr_cvar_t}
\end{equation}
where $\hat{\nu}$ is the MLE-estimated degrees-of-freedom, $t^{-1}(\alpha;\nu)$
is the Student-$t$ quantile function, and $f_t(\cdot;\nu)$ is its PDF.

\subsection{Filtered Historical Simulation (FHS)}

When \texttt{fhs=True}, the parametric distribution is replaced by the
empirical distribution of standardised residuals:
\begin{equation}
\hat{z}_t = \frac{r_t - \hat{\mu}}{\hat{\sigma}_t}
\end{equation}
VaR quantiles are then taken from the empirical distribution of $\{\hat{z}_t\}$:
\begin{equation}
q_\alpha^{\text{FHS}} = Q_{1-\alpha}\!\left(\{\hat{z}_t\}\right)
\end{equation}
(Note: the upper tail quantile $Q_{1-\alpha}$ is used to match the sign
convention for losses.)

FHS retains empirical tail behaviour while using GARCH-filtered volatility,
combining the strengths of both parametric and non-parametric approaches.

\subsection{Holding-Period Scaling}

\begin{equation}
\widehat{\text{VaR}}_\alpha^{(H)} = \widehat{\text{VaR}}_\alpha^{(1)} \cdot \sqrt{H}
\end{equation}

%-----------------------------------------------------------------------
\section{Parameters}

\begin{tabular}{llll}
\toprule
Parameter & Type & Default & Description \\
\midrule
\texttt{window}   & int   & 250   & Lookback window \\
\texttt{rescale}  & float & 1000  & Rescaling factor for returns (numerical stability) \\
\texttt{decay}    & float & 1     & EWMA pre-weighting ($\lambda=1$: no decay) \\
\texttt{ftol}     & float & 1e-6  & Optimizer convergence tolerance \\
\texttt{max\_iter}& int   & 150   & Maximum optimizer iterations \\
\texttt{fhs}      & bool  & False & Use Filtered Historical Simulation \\
\texttt{distr}    & str   & \texttt{'Normal'} & Innovation distribution: \texttt{'Normal'} or \texttt{'t'} \\
\texttt{qtls}     & list  & [0.05, 0.01, 0.001] & Tail probability levels \\
\texttt{holding\_period} & int & 1 & Forecast horizon in days \\
\bottomrule
\end{tabular}

%-----------------------------------------------------------------------
\section{Advantages and Limitations}

\begin{tabular}{p{6cm} p{8cm}}
\toprule
Advantages & Limitations \\
\midrule
Captures volatility clustering & Computationally heavier than EWMA/HS \\
Models leverage effect explicitly & Requires sufficient history for MLE convergence \\
Parametric tail estimation & Stationarity constraints may bind \\
FHS option for non-parametric tails & Square-root scaling is approximate \\
\bottomrule
\end{tabular}

%-----------------------------------------------------------------------
\section{Usage Example}

\begin{lstlisting}[language=Python]
from risk_analytics import risk_engines as re

var_gjr = re.portfolio_vaR(
    rets_orig  = rets,
    wgts       = weights,
    engine     = 'gjr',
    fmt_engine = {
        'window'         : 250,
        'rescale'        : 1000,
        'decay'          : 1,
        'ftol'           : 1e-6,
        'max_iter'       : 150,
        'fhs'            : False,
        'distr'          : 't',
        'qtls'           : [0.05, 0.01, 0.001],
        'holding_period' : 1,
    }
)
\end{lstlisting}

%-----------------------------------------------------------------------
\section{References}

Glosten, L.R., Jagannathan, R., and Runkle, D.E. (1993).
\textit{On the Relation between the Expected Value and the Volatility of
the Nominal Excess Return on Stocks.}
Journal of Finance, 48(5), 1779--1801.

\end{document}
